\documentclass{CU_IOPStyle_PXT999}
\DeclareUnicodeCharacter{22C6}{\ensuremath{\star}}

\usepackage{dblfloatfix}
\usepackage{amsmath}
\usepackage{subcaption}
\usepackage{placeins}
\usepackage{graphicx}
\graphicspath{{Figures/}}
\usepackage{cuted}
\usepackage{caption}
\usepackage{booktabs}
\usepackage{amssymb}
\usepackage{float}
\setlength{\fboxsep}{0pt}
\setlength{\fboxrule}{0.5pt}
\usepackage{longtable}
\usepackage{pdfpages}


\hypersetup{
    colorlinks=true,
    allcolors=blue
}


\begin{document}

\twocolumn[
\title{GaussCluster: A 3D Gaussian Mixture Model for Open Cluster Membership}

\author{Shakoor}{Abdul Raffey}{a}
\author{Roche}{Paul}{a}

\affiliation{a}{Department of Physics and Astronomy, Cardiff University, Cardiff CF24 3AA, United Kingdom}

\maketitle

\begin{abstract}
\textit{Context.} Isolating true dynamical members from the galactic background is essential for characterizing open clusters. In highly obscured environments like the Cygnus X region, spatial density algorithms struggle, often sweeping up heavy background contamination.

\textit{Aims.} We develop a two-stage methodology which couples blind spatial-density searches with a probabilistic kinematic extraction, enabling the recovery of a cluster's kinematic core even under extreme reddening conditions.

\textit{Methods.} We present \texttt{GaussCluster}, a 3D Gaussian Mixture Model pipeline utilizing high-precision Gaia DR3 proper motions and parallaxes. To evaluate the pipeline's efficiency, we applied it to the heavily obscured young cluster Berkeley 87 and a less complex control target, NGC 957. The derived kinematic populations were then cross-matched with theoretical PARSEC isochrones to establish their fundamental physical parameters.

\textit{Results.} We isolated 70 highly probable (P $\ge 0.9$) kinematic members in Berkeley 87, and 180 members in NGC 957, by ingesting the core isolated by the global spatial model, STORM, and probabilistically resolving the physical extent of the cluster within a 15-arcminute boundary. We derived a distance of $\sim 1.8$ kpc for Berkeley 87, alongside physical parameters for NGC 957 (10 Myr, 2.2 kpc), consistent with literature values. The 3D kinematic filtering revealed a photometric sequence yielding an isochrone age of $\sim$ 4.5 Myr for Berkeley 87, aligning with the $\le 5$ Myr physical limits dictated by the presence of post-main sequence stars.

\textit{Conclusions.} The purely kinematic, probabilistic workflow of \texttt{GaussCluster} provides an alternative approach to overcoming the observational biases caused by clumpy interstellar extinction and spatial overlapping. This methodology provides a foundation for building accurate catalogues of young stellar populations in complex environments.
\end{abstract}

\par\noindent\rule{\textwidth}{0.5pt}
\vspace{0.5cm}
]

\section{Introduction}
\label{ssec:1}

\subsection{Overview}
\label{ssec:1.1}
The accurate determination of open cluster parameters such as age, distance, and initial mass function is central to studies of galactic structure and stellar evolution \cite{portegies_zwart_star_2001}. However, extracting a clean catalogue of physically bound cluster members from the dense galactic background is a complex challenge, particularly in obscured regions of the Milky Way \cite{cantat-gaudin_milky_2022, de_la_fuente_clustered_2021}. This paper introduces \texttt{GaussCluster}, a 3D kinematic modelling pipeline designed to address this membership problem using high-precision data from the Gaia mission \cite{gaia_collaboration_gaia_2023}.

\subsection{Open Clusters as Astrophysical Laboratories}
\label{ssec:1.2}
Open clusters are widely recognised as valuable tools available for observational astronomy \cite{cantat-gaudin_milky_2022}. As gravitationally bound groups of stars born from the same giant molecular cloud at approximately the same epoch and distance, they provide astronomers with a natural laboratory for studying stellar evolution \cite{portegies_zwart_star_2001}, particularly in their early embedded phases \cite{lada_embedded_2003}. While theoretical isochrone fits to a cluster's Colour-Magnitude Diagram (CMD) are traditionally used to constrain age, distance, and reddening \cite{bressan_span_2012}, the advent of Gaia allows distance to be constrained kinematically, breaking photometric degeneracy. This makes open clusters valuable calibrators of stellar evolution timelines \cite{von_hippel_galactic_2005} and theoretical models \cite{bressan_span_2012}, particularly those involving rotational mixing \cite{meynet_stellar_2000}. Comparing cluster observations to theoretical tracks like the PARSEC models provides a direct test of stellar interior physics, like convective overshooting, rotation, and mass loss \cite{bressan_span_2012}.

Beyond the study of individual stars, the macroscopic spatial distribution and ages of open clusters across the galactic disc are useful in tracing spiral arm structures and understanding the broader history of star formation within the Milky Way \cite{dias_new_2002, vande_2010}. Historically, cluster classification helped map the Galaxy \cite{ruprecht_classification_1966, janes_galactic_1994}, while modern surveys continue to trace spiral arm structures with unprecedented precision \cite{cantat-gaudin_painting_2020}. 

Young open clusters, with ages on the order of a few to tens of millions of years, are important as they have not had sufficient time to dynamically dissolve into the galactic field \cite{lada_embedded_2003}. Because they have not yet been destroyed by two-body relaxation and tidal stripping, they retain a population of massive, early-type (O and B) stars \cite{portegies_zwart_star_2001}. These stars evolve rapidly and inject substantial amounts of mechanical and radiative energy into the ISM via stellar winds and subsequent core-collapse supernovae, making them sensitive tracers of recent star formation activity \cite{phelps_young_1993}. The initial mass function (IMF) of these young clusters also provides valuable insight into the star formation process itself, helping to answer questions about whether stellar mass distributions are universal across different galactic environments \cite{kroupa_variation_2001}.

\subsection{The Impact of High-Precision Astrometry}
\label{ssec:1.3}
Deriving accurate cluster parameters depends on the observer's ability to correctly separate physically bound cluster members from foreground and background field stars along the line of sight \cite{balaguer-nunez_determination_1998, sanders_71}. Before the onset of high precision astrometric surveys, membership was assigned photometrically \cite{janes_1988, gimenez_rgu_1980}. This involved plotting all stars in a given field on a CMD and statistically subtracting field stars based on a theoretical isochrone or empirical main sequence \cite{yadav_optical_2008}. Such methods have been standard practice for decades \cite{gimenez_rgu_1980}, although they are susceptible to foreground contamination \cite{gerasimenko_investigation_1991, Maciejewski_2007}. 

In crowded or reddened regions along the galactic plane, reddened foreground stars or luminous background objects can easily appear as highly obscured cluster members \cite{Maciejewski_2007, cantat-gaudin_milky_2022}. The accurate proper motions and parallaxes provided by Gaia Data Release 2 (DR2) \cite{lindegren_2018, luri_2018} and Data Release 3 (DR3) \cite{lindegren_2021} allow cluster membership to be assigned kinematically, largely independent of a cluster's photometric properties \cite{gaia_collaboration_gaia_2023}. 

This has started a revolution in cluster science, as astronomers have used Gaia kinematics to revise cluster parameters, discover entirely new clusters hidden in the galactic plane, and discard objects obtained from earlier photometric-only studies \cite{cantat-gaudin_gaia_2018}. Recent works including density-based spatial clustering algorithms like HDBSCAN \cite{hunt_improving_2023}, machine learning models such as Random Forests \cite{Mahmudunnobe_2021}, and probabilistic clustering models like UPMASK \cite{krone-martins_upmask_2014}, all use this data. Together, these projects have refined parameters for specific clusters \cite{bisht_deep_2021, koc_photometric_2022}, resolved complex star-forming regions \cite{de_la_fuente_clustered_2021}, and generated large scale, all-sky cluster catalogues \cite{cantat-gaudin_painting_2020}. However, massive homogeneous catalogues rely on global algorithmic settings. While these global settings excel at discovering average clusters across the whole sky, they struggle or fail entirely when confronted with peculiar, extreme differential reddening environments like the Cygnus X region \cite{hunt_improving_2023, de_la_fuente_clustered_2021}. 

\subsection{The Cygnus X Region and the ON2 Complex}
\label{ssec:1.4}
To fully understand the observational challenges associated with Berkeley 87 (hereafter: Be87), it is important to contextualise its location within the Milky Way. The cluster resides deep within the Cygnus X region, one of the most active star-forming regions in the local galactic neighbourhood \cite{reipurth_2008}. Spanning several degrees across the sky, Cygnus X contains millions of solar masses of molecular gas and dust, driving star formation that powers numerous HII regions and OB stars \cite{de_la_fuente_clustered_2021}. Within this area resides the ON2 star-forming region, characterised by variable interstellar extinction \cite{WEBDA}. The molecular clouds in this specific region absorb optical light, making observations at shorter optical wavelengths challenging \cite{Cardelli_1989}. 

Because the cluster is still physically present within, or lying immediately behind its natal molecular cloud, the dust distribution is clumpy \cite{lada_embedded_2003, de_la_fuente_clustered_2021}. This clumpy structure is what generates the severe differential reddening that affects all optical photometric studies of the cluster, as photons from adjacent stars must travel different column densities of interstellar material before reaching Earth \cite{de_la_fuente_clustered_2021}. 

\subsection{Be87: A Complex and Obscured Target}
\label{ssec:1.5}
\begin{figure}[htbp]
    \centering
    {\includegraphics[width=0.95\columnwidth]{Be87_Cluster_Image(2).png}}
    \caption{H$\alpha$ image of Be87 using the LCO 0.4m telescope, with clear dust obscurity visible}
    \label{fig:be87_image}
\end{figure}

The first complete observational characterisation of Be87 was carried out by Turner \& Forbes (1982) \cite{turner_berkeley_1982}, who identified it as host to several highly evolved stars, which includes the emission-line variable V439 Cyg and the oxygen-rich Wolf-Rayet star WR 142. The presence of WR 142 is diagnostic as it indicates that Be87 must have formed massive progenitor stars ($> 40 M_{\odot}$) \cite{sokal_chandra_2010}, placing constraints on the cluster's upper age limit, as such stars typically detonate as core-collapse supernovae within $\sim5$ million years \cite{crowther_2007}.

High-energy X-ray observations have also provided a characterisation of this stellar population. Observations by the Chandra X-ray Observatory confirmed the membership of WR 142 and several OB-type stars while revealing hot thermal plasma emission generated by colliding stellar winds, typical of dense, young clusters \cite{sokal_chandra_2010}, providing further evidence of the cluster's extreme environment \cite{skinner_chandra_2019}. 

More recently, a study by de la Fuente et al. (2021) \cite{de_la_fuente_clustered_2021} compiled a multi wavelength survey of the Be87/ON2 field, providing a Gaia EDR3-based distance determination alongside a catalogue of confirmed cluster members with known spectral types. This work addressed the difficulty of handling overlapping young stellar populations along this crowded line of sight \cite{de_la_fuente_clustered_2021}, a task which previous automated spatial algorithms had struggled to map accurately due to the clumpy dust distribution \cite{hunt_improving_2023}. 

This project aims to directly address this spatial-overlapping limitation, by establishing fundamental parameters through Gaussian kinematic membership selection. Evaluating highly obscured clusters requires confidence in the methodology, so this study incorporates an independent check on a separate open cluster, NGC 957. By comparing this purely kinematic model against a spatial density-based algorithm on the control target (NGC 957), this study aims to demonstrate how a Gaussian mixture model can filter out spatial contamination before deploying the pipeline to resolve the highly obscured Be87.

The structure of this paper is as follows: Section \ref{ssec:2} outlines the foundational methodology and astrometric data acquisition. Section \ref{ssec:3} details the architecture of the \texttt{GaussCluster} pipeline, Section \ref{ssec:4} presents the kinematic recovery of Be87 and NGC 957, and details the derivation of parameters via isochrone fitting. Section \ref{ssec:5} provides a summary of the findings, and Section \ref{ssec:6} concludes with a brief discussion of future works.

\section{Methodology}
\label{ssec:2}
The methodology implemented in this study is centred on the utilization of precision astrometry from the Gaia mission to isolate true cluster members from the galactic background.

\subsection{Astrometric Data Acquisition and Filtering}
\label{ssec:2.1}
To determine the kinematic structure of Be87, astrometric data was retrieved from the third data release of the Gaia mission (Gaia DR3) \cite{gaia_collaboration_gaia_2023} via the \texttt{astroquery} Python package \cite{ginsburg_astroquery_2019}. A circular cone search was performed centred on the accepted cluster coordinates ($\alpha = 20^\text{h} 21^\text{m} 35.0^\text{s}$, $\delta = +37^\circ 25' 59''$, J2000)\cite{Tarricq_2021}. The search radius was restricted to 15 arcminutes. This boundary was selected to capture the inner core and primary halo of the cluster while mitigating the field contamination inherent to the Cygnus X region. This spatial restriction is consistent with the typical structural extents and tidal radii of open clusters residing within the galactic disc \cite{Tadross01062014}. While stars in the extended outer halo may have been ejected from the core via dynamical interactions \cite{spitzer_1987, baumgardt_2003, oh_kroupa_2016}, limiting the radius to 15 arcminutes ensures the subsequent kinematic models are not overwhelmed by the non-uniform background field.

To ensure high astrometric accuracy before applying any statistical models, the initial dataset was filtered. The Gaia catalogue contains $\sim 27000$ sources in this region, but many suffer from poor astrometric solutions \cite{gaia_collaboration_gaia_2023, lindegren_2021}. Only targets where the measurement errors in proper motion ($\epsilon_{\mu_\alpha *}$, $\epsilon_{\mu_\delta}$) were small ($\le 0.5$ mas/yr) were retained. The catalogue is then filtered based on the Renormalised Unit Weight Error (RUWE). A high RUWE typically indicates a poor fit to the single-star astrometric model, often caused by unresolved binary companions perturbing the photocenter \cite{belokurov_2008}. A RUWE $< 1.4$ cutoff was applied, to ensure the astrometric solutions were consistent \cite{lindegren_2018}. This pre-processing step removed inaccurate measurements, unresolved close binaries, and faint background stars with uncertain kinematics that could artificially inflate the calculated cluster velocity dispersions \cite{gieles_2010}.

\subsection{Kinematic Membership Determination}
\label{ssec:2.2}
A limitation of traditional spatial and photometric studies of Be87 is the assumption that stars clumped together, visually, in the sky belong to the cluster \cite{de_la_fuente_clustered_2021}. To isolate dynamical members from the dense Cygnus X foreground and background, a probabilistic approach was adopted, expanding upon the bivariate mixture model framework initially formulated by Balaguer-N\'u\~nez et al. (1998) \cite{balaguer-nunez_determination_1998}. While the traditional model relies purely on 2D proper motions, this approach is susceptible to line-of-sight contamination. Therefore, the model was adapted into a 3D Gaussian Mixture Model (3GMM) by incorporating Gaia parallax ($\varpi$) as a third kinematic dimension, extending the model into a fully three-dimensional spatial-kinematic framework.

In this expanded framework, the distribution of all stars is modelled as a multivariate normal distribution defined by a data vector $\mathbf{x} = (\mu_{\alpha*}, \mu_{\delta}, \varpi)^T$. For a given star with corresponding observational covariance $\Sigma_{\epsilon}$, the probability density function for the cluster, $\Phi_c$, is generalized as \cite{zhao_1990}:
\begin{equation}
    \resizebox{0.95\columnwidth}{!}{$\displaystyle
        \Phi_c = \frac{1}{(2\pi)^{3/2} |\Sigma_c + \Sigma_{\epsilon}|^{1/2}} \exp \left\{ -\frac{1}{2} (\mathbf{x} - \boldsymbol{\mu}_c)^T (\Sigma_c + \Sigma_{\epsilon})^{-1} (\mathbf{x} - \boldsymbol{\mu}_c) \right\}
    $}
\end{equation}
where $\boldsymbol{\mu}_c$ represents the systemic 3D centroid of the cluster (proper motions and parallax) and $\Sigma_c$ is its intrinsic dispersion matrix.

However, because the background field stars in the Cygnus plane are highly non-uniform and cannot be described by a single broad Gaussian, the field distribution ($\Phi_f$) was modelled using a multi-dimensional Kernel Density Estimate (KDE) \cite{cabrera-cano_1990, krone-martins_upmask_2014}. By treating the spatial-kinematic field data as a sum of local Gaussian kernels, the KDE models the non-Gaussian structure of the field. For a given data point $\mathbf{x}$, the field probability density is evaluated as:

\begin{equation}
    \Phi_f(\mathbf{x}) = \frac{1}{N_f} \sum_{i=1}^{N_f} K_{\mathbf{H}} (\mathbf{x} - \mathbf{x}_i)
\end{equation}

where $N_f$ is the number of field stars used to train the KDE, $K_{\mathbf{H}}$ is the multivariate Gaussian kernel, and $\mathbf{H}$ is the bandwidth matrix determining the smoothing scale (calculated via Scott's rule \cite{scott_1992}).
The final kinematic membership probability $P$ for each star is then defined as the ratio of the cluster density to the total combined density (cluster plus field)\cite{vasilevskis_1958, sanders_71}:

\begin{equation}
    P = \frac{n_c \cdot \Phi_c}{n_c \cdot \Phi_c + n_f \cdot \Phi_f}
\end{equation}

where $n_c$ and $n_f$ represent the normalized fractions of cluster and field stars in the region, respectively ($n_c + n_f = 1$). These parameters were solved using the empirical Gaia data, giving a converged catalogue of 3D kinematic members. A probability threshold of $P \ge 0.9$ was adopted, for the derivation of parameters.


\begin{figure*}[t]
    \centering
    \includegraphics[width=0.95\textwidth]{GaussCluster_Diagnostics_Cosmetic_Be87.png}
    \caption{\texttt{GaussCluster} 4-panel diagnostic dashboard for Be87. The algorithm recovered 70 members ( $\ge 0.9$), forming a tight Main Sequence.}
    \label{fig:be87_diagnostics}
\end{figure*}

\begin{figure*}[t]
    \centering
    \includegraphics[width=0.95\textwidth]{Be87_GaussCluster_vs_STORM.png}
    \caption{Spatial and parallax distribution comparison, along with a Colour-Magnitude Diagram for Be87. \texttt{GaussCluster} model (blue, 70 members) versus STORM (red, 146 members within 15 arcminutes).}
    \label{fig:storm_comparison_be87}
\end{figure*}

\subsection{Photometric Validation via Gaia DR3}
\label{ssec:2.3}
To validate the kinematic membership model, a CMD was constructed, using the native Gaia DR3 passbands ($G$, $G_{BP}$, $G_{RP}$) for all stars identified as highly probable members ($P \ge 0.9$). Rather than relying on external ground-based photometric catalogues, utilising Gaia photometry allows homogeneity \cite{riello_2021}, allowing a direct visual confirmation that the kinematically selected stars form a distinct, physical main sequence. 

Theoretical stellar evolutionary tracks were generated from the PARSEC database \cite{bressan_span_2012} assuming solar metallicity ($Z=0.0152$), which is appropriate for young galactic disc objects \cite{netopil_2016}. Rather than simultaneously attempting to fit age, distance, and reddening, an approach which frequently leads to degenerate parameter solutions \cite{anders_2004}, the distance modulus was anchored using the median zero-point-corrected Gaia DR3 parallax of the high-probability kinematic members \cite{lindegren_2021}, following the statistical parallax methodologies outlined by \cite{luri_2018}. This calculation allowed us to treat distance as a fixed prior. By locking the vertical position of the isochrone, the isochrone's photometric properties (age, reddening) could then be compared against the kinematically selected population without additional free parameters \cite{meynet_stellar_2000}.

\begin{figure*}[t]
      \centering
      \includegraphics[width=0.95\textwidth]{NGC957_GaussCluster_vs_STORM.png}
      \caption{Spatial and parallax distribution comparison, along with a Colour-Magnitude Diagram for NGC 957. \texttt{GaussCluster} model (blue, 180 members) versus STORM (red, 311 members within 15 arcminutes).}
      \label{fig:ngc957_storm}
\end{figure*}


\section{The GaussCluster Pipeline}
\label{ssec:3}
Traditional open cluster studies often rely on manual, tailored code that is calibrated to a single target \cite{cantat-gaudin_milky_2022, hunt_improving_2023}. Rather than writing target-specific analytical scripts, the 3D kinematic framework was engineered into a reusable Python Command Line Interface (CLI) tool, \texttt{GaussCluster} \footnote{The source code for the \texttt{GaussCluster} pipeline is publicly available at \url{https://github.com/raffey05-sudo/GaussCluster}}.

\texttt{GaussCluster} was designed to automate the process of separating intrinsic cluster kinematics from background field contamination. The architecture of the software is built around a sequential pipeline:

\begin{enumerate}
    \item \textbf{Data Acquisition \& Filtering:} In this integrated pipeline, the output generated by STORM \cite{storm2026} is ingested as initial 'seed' members for the target cluster, into \texttt{GaussCluster} \cite{gausscluster2026}. The software interfaces directly with the Gaia DR3 archive via the \texttt{astroquery} API \cite{ginsburg_astroquery_2019} to download the complete stellar field out to a customizable search radius, which the user defines based on the expected tidal extent of their specific target. To prevent astrometric noise or unresolved binary systems from artificially inflating the velocity dispersion, the software automatically purges any source with a RUWE $\ge 1.4$ \cite{lindegren_2018} or proper motion errors exceeding 0.5 mas/yr.
    
    \item \textbf{3D Kinematic Core Anchoring:} Unlike blind algorithms \cite{storm2026}, \texttt{GaussCluster} anchors its first step by using a 3D sigma-clipping algorithm on the proper motions ($\mu_{\alpha*}, \mu_{\delta}$) and parallax ($\varpi$) of the ingested core members to calculate an initial estimate for the cluster's centroid. The cluster distribution is then modelled as a trivariate Gaussian.
      
    \item \textbf{Non-Parametric Field KDE:} The software employs a Non-Parametric Kernel Density Estimator (KDE) to map the asymmetric astrometric density of the surrounding galactic field. This non-parametric approach to field modelling has become standard practice in probabilistic membership selection \cite{krone-martins_upmask_2014}.
    
    \item \textbf{Dynamic Bayesian Prior:} The algorithm evaluates the kinematic density of stars in the core and subtracts the expected KDE field density, dynamically calculating the local cluster density to establish a Bayesian prior ($n_c$). This helps reduce manual steering and allows the software to automatically discover the full physical scale of the cluster.
    
    \item \textbf{Diagnostic Export:} Finally, the software extracts all stars above a threshold ($P \ge 0.9$) and generates a diagnostic dashboard, based strictly on the final 3GMM kinematic probabilities. This can be seen in Figure \ref{fig:be87_diagnostics} which shows the output of \texttt{GaussCluster} when run on Be87.
\end{enumerate}


\subsection{Methodological Choices and Limitations}
\label{ssec:3.1}
While the 3D kinematic Gaussian Mixture Model isolates cluster members in heavily obscured regions, the implementation of the algorithm required multiple methodological decisions to ensure stability and computational efficiency. 

\subsubsection{Decoupled Field Contamination Modeling} 
\label{ssec:3.1.2} 
The field-star mixing fraction, $n_f$, was estimated independently of the KDE field density model. A geometric argument based on the excess source density beyond the core exclusion radius was used, assuming approximate local uniformity of field sources in normalised kinematic space \cite{cantat-gaudin_gaia_2018, Sarro_2014}. This separates the total number of expected background stars from the actual shape of their astrometric distribution, meaning the two are not required to be mutually consistent. This two-stage decoupling approach is often utilised in astrometric cluster models to prevent local kinematic anomalies from artificially inflating the expected global field density \cite{dias_2014}.

\subsubsection{Mahalanobis Distance Truncation}
\label{ssec:3.1.1}
To ensure that outliers do not artificially widen the cluster model during the iterative loop, the membership probability for any source lying beyond a Mahalanobis distance of $3\sigma$ was forced to zero \cite{rosseeuw_1190}. This is equivalent to fitting a shortened, rather than a full, trivariate Gaussian to the astrometric distribution at each Expectation-Maximization (E-M) step, introducing a downward bias in the recovered velocity and parallax dispersion relative to an untruncated fit, as high-dispersion tail members are excluded \cite{Huber_1981, johnson_1994}. Standard E-M algorithms are sensitive to outlier contamination, a limitation of Finite Mixture Models \cite{mclachlan_2000}, so a $3\sigma$ truncation was applied \cite{krone-martins_upmask_2014}.

\subsubsection{Static KDE Background Estimation} 
\label{ssec:3.1.3} 
For algorithmic stability, the background field KDE was fit only once using the field sample defined by the initial core-based cluster estimate, rather than being continuously refit at each E-M iteration \cite{cabrera-cano_1990, dias_2014}. This implementation assumes that the core-field boundary does not shift enough during refinement to alter the field sample's composition. Utilizing a static field distribution from a surrounding region is used to prevent iterative instability during maximum likelihood estimation \cite{zhao_1990}. Given the stable convergence observed, this assumption is well supported. The kinematically distinct nature of the initial STORM seed cores ensures that the initial field selection remains representative throughout the model's execution.


\section{Results and Discussion}
\label{ssec:4}

The standard output generated by the \texttt{GaussCluster} pipeline is a 4-panel diagnostic dashboard, visualized in Figure \ref{fig:be87_diagnostics} for Be87. Reading from left to right, the dashboard displays the isolated cluster's proper motion Vector Point Diagram (VPD), its spatial footprint (RA/Dec), its parallax distribution, and the resulting Colour-Magnitude Diagram (CMD). To evaluate the two-stage pipeline, this expanded \texttt{GaussCluster} output was analyzed in the context of the initial STORM core input and the recent catalogue published by Hunt (2023) \cite{hunt_improving_2023}.

\subsection{Kinematic Structure of Berkeley 87}
\label{ssec:4.1}

Using the STORM extraction as the initial seed core, the application of the 3GMM yielded 70 highly probable members for Be87 at the $P \ge 0.9$ confidence threshold. The complete catalogue of these 70 core kinematic members, including their astrometric parameters and final membership probabilities, is provided in Appendix \ref{app:catalogue}.

The difference in total membership size from STORM (146 members within 15 arcminutes) to GaussCluster (70 members) highlights the divergence between spatial-density and kinematic approaches. Density-based algorithms like HDBSCAN (the clustering engine behind STORM and Hunt 2023) rely on identifying continuous regions of high point density. In crowded galactic plane fields like Cygnus X, the dense background can bridge with the cluster, causing density algorithms to include larger numbers of field stars.

By operating in three-dimensional kinematic space ($\mu_{\alpha*}$, $\mu_{\delta}$, $\varpi$) and utilizing the full observational covariance matrix, \texttt{GaussCluster} applies a stricter kinematic constraint. It rejects spatial interlopers that do not share the cluster's precise systemic motion. The inclusion of parallax as a third kinematic dimension ensures that even if a star resides within the spatial core, it is discarded if it does not share the same distance modulus. This demonstrates that while density models encompass a wider spatial footprint, probabilistic 3D kinematic models return a more conservative, kinematically restricted physical population. To verify that this restricted core reflects a gravitationally bound structure, empirical King Profile models \cite{king_structure_1962} were fitted to the spatial distributions of both clusters (see Appendix \ref{app:king}).

Note that when comparing the resultant parallax distributions, as seen in both Figure \ref{fig:storm_comparison_be87} and Figure \ref{fig:ngc957_storm}, a slight systemic offset is visible between \texttt{GaussCluster} and the literature catalogues. This occurs because the \texttt{GaussCluster} pipeline explicitly applies the Lindegren et al. (2021) \cite{lindegren_2021} zero-point corrections to the native Gaia DR3 parallaxes, whereas the comparison catalogues rely on uncorrected values.

\begin{figure}[h]
    \centering
    \includegraphics[width=0.95\columnwidth]{NGC957_CMD_Isochrone.png}
    \caption{Gaia DR3 Colour-Magnitude Diagram for the kinematically selected members of NGC 957. The PARSEC isochrone fit anchors the distance modulus, with the pre-main sequence transition zone visible at the faint end.}
    \label{fig:ngc957_cmd}
\end{figure}

\subsection{Methodology Validation: NGC 957}
\label{ssec:4.2}

To ensure the pipeline was not biased by the spatial density and structural complexity of the Cygnus X field, an independent validation of \texttt{GaussCluster} was performed on the open cluster NGC 957. Ingesting the STORM core (311 members) as the initial seed, \texttt{GaussCluster} extracted a final 3D population of 180 members, where 165 of the 180 members (91.7\%) are also STORM members. Cross-matching this kinematic output with the independent \cite{hunt_improving_2023} catalogue gave an intersection of 134 common members.

Both the STORM and Hunt methodologies extracted significantly more total members for NGC 957 (311 and 252 respectively) than \texttt{GaussCluster}. Consistent with the results from Be87, these variations are attributable to clustering mechanics. Because both STORM and Hunt rely on spatial-density algorithms, they are prone to absorbing faint background sources that spatially overlap with the cluster. The 3GMM mathematically filters out these spatial interlopers, demonstrating that the pipeline consistently recovers a more kinematically restricted physical sequence, regardless of the cluster's environment.

\begin{figure}[!h]
    \centering
    \includegraphics[width=0.95\columnwidth]{Be87_CMD_Isochrone.png}
    \caption{Gaia DR3 Colour-Magnitude Diagram for the kinematically selected members of Be87, demonstrating the isolation of a distinct main sequence population against the extreme differential reddening of Cygnus X.}
    \label{fig:be87_cmd}
\end{figure}

\subsection{Photometric Validation and Fundamental Parameters}
\label{ssec:4.3}

\begin{table}[h]
    \centering
    \caption{Fundamental parameters derived for NGC 957 and Be87 compared against literature values.}
    \vspace{0.2cm}
    \begin{tabular}{lcc}
        \hline\hline
        Parameter & NGC 957 & Berkeley 87 \\
        \hline
        \multicolumn{3}{c}{\textit{This Work}} \\
        \hline
        Age (Myr) & 10.0 & 4.47 \\
        Distance (kpc) & 2.2 & 1.78 \\
        \hline
        \multicolumn{3}{c}{\textit{Literature Mean}} \\
        \hline
        Age (Myr) & 10.0\textsuperscript{a} & 2.7\textsuperscript{b} \\
        Distance (kpc) & 2.0\textsuperscript{a} & 1.45\textsuperscript{b} \\
        \hline\hline
    \end{tabular}
    \label{tab:fundamental_parameters}
    \vspace{1ex}
    
    {\footnotesize 
    \textsuperscript{a}Mean of \cite{yadav_optical_2008} (10 Myr, 2.2 kpc) and \cite{cantat-gaudin_painting_2020} (10 Myr, 1.8 kpc).\\
    \textsuperscript{b}Mean of \cite{turner_berkeley_1982} (1-2 Myr, 1.2 kpc) and \cite{cantat-gaudin_painting_2020} (4 Myr, 1.7 kpc).}
\end{table}

To validate the kinematic isolation, CMDs were constructed using the native Gaia DR3 passbands. Theoretical stellar evolutionary tracks were generated from the PARSEC database \cite{bressan_span_2012} assuming a single solar metallicity ($Z=0.0152$) for both clusters. The distance modulus for each cluster was mathematically anchored using the median zero-point-corrected Gaia DR3 parallax of the members \cite{lindegren_2021}. Rather than relying on ages from the literature, the literature values were used to constrain the bounds of the theoretical isochrone generation. For Be87, tracks were generated between $\log(\text{Age}) = 6.30$ and $6.70$ (2 to 5 Myr), while for NGC 957 they were generated between $6.60$ and $7.05$ (4 to 11 Myr), both utilizing a 0.05 logarithmic step. For NGC 957, the derived distance of 2.2 kpc is slightly larger than the Gaia-based literature average (1.8 kpc), a discrepancy which may arise from the specific choice of PARSEC isochrone metallicity.

A prominent widening is visible at the faint end of the main sequence ($G \gtrsim 15$) for both clusters. This feature is consistent with the Pre-Main Sequence (PMS) transition zone, though unresolved binaries and differential reddening also contribute to this spread \cite{cantat-gaudin_gaia_2018, cantat-gaudin_milky_2022}. At ages of $\sim10$ Myr or less, low-mass stars have not yet fully contracted onto the Zero-Age Main Sequence (ZAMS) \cite{Soderblom_2014}. The PARSEC evolutionary tracks explicitly model this delayed contraction, which matches the rightward hook toward cooler temperatures observed in the Gaia photometry.

\subsubsection{Isochrone Ages}
\label{ssec:4.3.1}

The theoretical isochrone fit for NGC 957 yields an age of 10 Myr ($\log(\text{Age}) = 7.0$), which matches the literature mean (Table \ref{tab:fundamental_parameters}). For Be87, the updated kinematic isolation and clean main sequence allowed for an isochrone fit yielding an age of roughly 4.5 Myr ($\log(\text{Age}) = 6.65$). Decades of multi-wavelength observations establish Be87 as an extremely young embedded cluster, with literature ages ranging from 1 to 5 Myr and distances from 1.2 to 1.7 kpc \cite{turner_berkeley_1982, cantat-gaudin_painting_2020}. 

While our derived age (4.47 Myr) and distance (1.78 kpc) differ from the historical literature mean, they remain consistent with the $\le 5$ Myr physical limit dictated by the cluster's massive post-main sequence population \cite{sokal_chandra_2010, crowther_2007}. The 3D kinematic filtering returned a clean photometric sequence that supports these young age limits.

Furthermore, an independent photometric survey utilizing the LCO 0.4m network placed H$\alpha$ emission constraints on the cluster. A complete lack of active circumstellar decretion disks was seen among the members in the FOV, including the known emission variable V439 Cyg, which restricts the current evolutionary state of the B-type population (Appendix \ref{app:halpha}).

\section{Conclusion}
\label{ssec:5}

The primary objective of this work was to develop and validate a three-dimensional kinematic pipeline, \texttt{GaussCluster}, capable of isolating open clusters from crowded galactic environments. By operating strictly in a continuous probabilistic space ($\mu_{\alpha*}$, $\mu_{\delta}$, $\varpi$) via a Gaussian Mixture Model, the pipeline provides an alternative to density-based clustering algorithms, such as HDBSCAN, which rely on spatial density thresholds. 

The software was validated on the open cluster NGC 957 and applied to the heavily obscured Cygnus X region to recover the Be87 cluster. The quantitative recovery metrics are summarized in Table \ref{tab:membership_overlap}.

\begin{table}[!h]
    \centering
    \caption{Membership overlap counts for Be87 and NGC 957 across the GaussCluster, STORM, and \cite{hunt_improving_2023} catalogues, strictly restricted to a 15-arcminute radius for a homogeneous spatial comparison. Percentages express each intersection as a fraction of \texttt{GaussCluster}'s final output size for that cluster.}
    \vspace{0.2cm}
    \begin{tabular}{lcc}
        \hline\hline
        Catalogue / Pipeline & Berkeley 87 & NGC 957 \\
        \hline
        GaussCluster (3GMM) & 70 & 180 \\
        STORM (15') & 146 & 311 \\
        Hunt (2023) (15') & 82 & 252 \\
        \hline
        Intersection (GC $\cap$ STORM) & 88.6\% & 91.7\% \\
        Intersection (GC $\cap$ Hunt) & 48.6\% & 74.4\% \\
        \hline\hline
    \end{tabular}
    \label{tab:membership_overlap}
\end{table}

The overlap analysis reveals differences in membership allocation depending on the clustering architecture. The strong overlap with STORM ($\sim 90$\%) confirms the pipeline reliably anchors onto the initial dense core, while the moderate intersection (74.4\%) with the \cite{hunt_improving_2023} catalogue for NGC 957 highlights differences in spatial and kinematic limitations. While Hunt's all-sky HDBSCAN approach isolates a larger population, the 3GMM's mathematical approach enables the probabilistic recovery of a more kinematically restricted core. Comprehensive visual dashboards comparing our model against the catalogue of Hunt et al. (2023) for Berkeley 87, as well as against both STORM and Hunt et al. (2023) for the validation cluster NGC 957, are available in Appendices \ref{app:be87_hunt} and \ref{app:ngc957}.

Photometric validation using Gaia DR3 passbands yielded distinct main sequences for both isolated kinematic groups. The isochrone fit for NGC 957 aligned strongly with the literature mean age of 10 Myr. Crucially, the fit for Be87 produced a young age of 4.5 Myr, which is consistent with the established 1--5 Myr age derived from decades of multi-wavelength observations \cite{turner_berkeley_1982, cantat-gaudin_painting_2020}. 

In summary, the \texttt{GaussCluster} pipeline provides a strict kinematic methodology for isolating cluster cores, while the photometric analysis demonstrates the ongoing need for complementary infrared observations or differential-reddening corrections when characterizing extremely young populations in the galactic plane.

\section{Future Work}
\label{ssec:6}
Future software development will focus on automating this two-stage process into a single, unified Graphical User Interface (GUI) or API wrapper. In this future release, users will simply specify a sky coordinate, allowing the software to trigger a STORM extraction (Stage 1) and pipe the output directly into \texttt{GaussCluster} (Stage 2) without requiring manual data-handling between the Python scripts. 

Future observational work must also focus on obtaining deep near-infrared photometry (e.g., via the UKIRT Infrared Deep Sky Survey or JWST) and high-resolution spectroscopic follow-up for these newly confirmed kinematic members. Such observations are needed to characterise embedded populations \cite{perren_analysis_2022}, determine precise cluster ages \cite{angelo_exploring_2025}, and assess the longevity of circumstellar disks \cite{alvarez-baena_longevity_2024}. 

By combining the blind discovery power of density searches with the kinematic precision of a 3GMM, this integrated tool will provide a framework for charting galactic open clusters.

\section*{Acknowledgements}
I would like to express my sincere gratitude to my supervisor, Prof. Paul Roche, for his practical guidance on cluster analysis and support during the development of this methodology. His expertise and encouragement were valuable to this project.

The computational framework and \texttt{GaussCluster} pipeline were built in Python, using its open-source scientific ecosystem. This includes Astropy \cite{astropy_2013}, NumPy \cite{Harris_2020}, SciPy \cite{scipy}, pandas \cite{pandas_2020}, scikit-learn \cite{scikit-learn}, and Matplotlib \cite{matplotlib}. 

Additionally, this research utilized the VizieR catalogue access tool and the SIMBAD astronomical database (operated at CDS, Strasbourg, France), alongside PARSEC theoretical stellar evolution models. Data retrieval used the Astroquery API \cite{ginsburg_astroquery_2019}.

Finally, this work has made use of data from the European Space Agency (ESA) mission \textit{Gaia} (\url{https://www.cosmos.esa.int/gaia}), processed by the \textit{Gaia} Data Processing and Analysis Consortium (DPAC, \url{https://www.cosmos.esa.int/web/gaia/dpac/consortium}). Funding for the DPAC has been provided by national institutions, in particular the institutions participating in the \textit{Gaia} Multilateral Agreement.

\section*{Data Availability}
The \texttt{GaussCluster} Python package developed for this dissertation is open-source and publicly hosted at \url{https://github.com/raffey05-sudo/GaussCluster}. Furthermore, the complete reproducibility repository for this manuscript containing the LaTeX source code, raw data catalogues, and all analytical scripts, is available at \url{https://github.com/raffey05-sudo/MSC-Dissertation}.


\clearpage
\bibliography{References} 
\newpage
\onecolumn
\include{appendix}

\clearpage
\newpage
\includepdf[pages=1-3]{SOC}
\end{document}
